% ====================================================================
% Mapa tematico de literatura y brecha de investigacion
% ====================================================================
\begin{figure}[H]
\centering
\resizebox{\textwidth}{!}{%
\begin{tikzpicture}[
  font=\small,
  line/.style={-{Latex[length=2.2mm]},thick,color=VIUOrange!90!black},
  stream/.style={draw=VIUDark!70,fill=VIULight,rounded corners=3pt,
    align=center,minimum width=3.35cm,minimum height=1.0cm,text width=3.15cm},
  note/.style={draw=VIUGray!55,fill=VIUWhite,rounded corners=3pt,
    align=left,text width=3.55cm,inner sep=5pt,font=\scriptsize},
  gap/.style={draw=VIUOrange,fill=VIUOrange!10,rounded corners=4pt,
    align=center,text width=5.35cm,inner sep=7pt,font=\small\bfseries},
  claim/.style={draw=VIUDark,fill=VIUDark!5,rounded corners=3pt,
    align=center,text width=4.65cm,inner sep=6pt,font=\scriptsize}
]
\node[stream] (wh) at (0,3.0) {AMR en almacenes\\\citep{leanh2006agv,wurman2008coordinating,azadeh2019warehouse,keith2024amrwarehouse}};
\node[note, right=0.6cm of wh] (whn) {Optimiza flujo, despacho, rutas y decisiones de sistema. Normalmente la carga movida por cada robot ya es compatible con su plataforma.};

\node[stream] (mrta) at (0,1.35) {MRTA, subastas y coaliciones\\\citep{gerkey2004formal,dias2006market,korsah2013taxonomy,choi2009consensus,dutta2021hedonic,aziz2021complexity}};
\node[note, right=0.6cm of mrta] (mrtan) {Formaliza asignaciones ST/MT, SR/MR e IA/TA. Resuelve elecciones discretas, a menudo separadas del control continuo.};

\node[stream] (pop) at (0,-0.3) {Juegos poblacionales y consenso\\\citep{sandholm2010population,quijano2017population,barreiro2017distributed,kia2019tutorial,martinezpiazuelo2022tcss}};
\node[note, right=0.6cm of pop] (popn) {Aporta leyes distribuidas, potenciales y convergencia. Suele trabajar con cuotas o formaciones fijadas, no con reclutamiento por demanda minima.};

\node[stream] (transport) at (0,-1.95) {Transporte cooperativo\\\citep{an2023cooperative,ebel2024cooperative,rosenfelder2024force,shan2024distributed}};
\node[note, right=0.6cm of transport] (trn) {Estudia movimiento conjunto y control de carga. La composicion de la coalicion suele estar dada o resuelta por otra capa.};

\node[gap, right=1.25cm of popn] (gap) {Brecha del TFM\\Una misma ley local debe estimar oferta, formar coaliciones por cota inferior, revisar estrategias y generar movimiento diferencial reactivo.};
\node[claim, right=0.9cm of gap] (tfm) {Aporte propuesto\\Smith + DAC + payoff de deficit + tasa espacial + campo reactivo, con prueba en caso homogeneo y extension por capacidad efectiva.};

\draw[line] (whn.east) -- (gap.west);
\draw[line] (mrtan.east) -- (gap.west);
\draw[line] (popn.east) -- (gap.west);
\draw[line] (trn.east) -- (gap.west);
\draw[line] (gap.east) -- (tfm.west);

\node[draw=VIUGray!60,fill=VIUWhite,rounded corners=3pt,align=center,
  font=\scriptsize,text width=8.8cm,below=0.6cm of gap] (risk)
  {Riesgo metodologico: si se afirma todo como garantia cerrada, el trabajo queda sobredimensionado. Por eso se separan nucleo probado, implementacion perturbada y extensiones futuras.};
\draw[line,VIUGray!70] (gap.south) -- (risk.north);
\end{tikzpicture}}
\caption[Mapa de literatura y brecha de investigación]{Mapa temático de literatura y brecha de investigación. Las familias existentes cubren partes del problema, pero no cierran simultáneamente la asignación por cota inferior, la revisión poblacional distribuida, el movimiento diferencial y la transición a transporte cooperativo.}
\label{fig:gap_literatura}
\end{figure}
